% Adapted from "Behavioral biology diagram" by Izaak Neutelings.
% https://tikz.net/behavior_sapolsky/
% Original and this derivative: CC BY-SA 4.0.
% Inspired by Robert Sapolsky's cross-timescale framework for behavior.
% Changes: semantic styles, neutral typography, spacing, and label placement.
\documentclass[tikz,border=8pt]{standalone}
\usepackage{minimalist-profile}
\usetikzlibrary{arrows.meta,backgrounds,decorations.pathreplacing}

\tikzset{
  bucket/.style={rounded corners=\msSpaceUnit,line width=\msLineContext,draw opacity=\msOpacityForeground,fill=none},
  branch/.style={draw=msColor5,ms structure},
  timespan/.style={decorate,decoration={brace,amplitude=\msSpaceLabel},ms structure}
}

\begin{document}
\begin{tikzpicture}[ms base,x=1.2cm,y=1.1cm]
  % Buckets precede marks on the main layer so the paper background cannot cover them.
    \draw[bucket,draw=msColor3]
      (-0.65,-2.35) rectangle (11.75,2.45);
    \draw[bucket,draw=msColor5]
      (-0.45,0.40) rectangle (4.75,1.90);
    \draw[bucket,draw=msColor1]
      (2.95,-1.80) rectangle (7.25,1.78);
    \draw[bucket,draw=msColor2]
      (6.00,0.38) rectangle (8.35,1.85);
    \draw[bucket,draw=msColor6]
      (7.35,-1.70) rectangle (9.90,1.75);


  \node[ms body,text=msColor5,anchor=south west] at (-0.35,1.92) {evolution};
  \node[ms body,text=msColor1,anchor=north west,yshift=-\msSpaceLabel] at (2.95,-1.80) {epigenetics};
  \node[ms body,anchor=south west] at (6.10,1.88) {\textcolor{msColor2}{\rule{\msMarkerPoint}{\msLineEmphasis}}\enspace endocrinology};
  \node[ms body,text=msColor6,anchor=south east] at (9.80,1.78) {neurobiology};
  \node[ms body,anchor=south east] at (11.55,-2.25) {\textcolor{msColor3}{\rule{\msMarkerPoint}{\msLineEmphasis}}\enspace ethology};

  \draw[ms arrow,draw=msInk,ms emphasis]
    (0,0) -- (10.25,0) node[right,ms title,text=msInk] {behavior};

  \draw[branch] (0.90,0.06) to[out=180,in=-90] ++(-0.48,0.48)
    node[above,ms body,text=msColor5] {species};
  \draw[branch] (2.40,0.06) to[out=180,in=-90] ++(-0.48,0.48)
    node[above,ms body,text=msColor5] {population};
  \draw[branch,draw=msInk] (4.20,0.06) to[out=180,in=-90] ++(-0.50,0.60)
    node[above,align=center,ms body] {individual\\genetics};
  \draw[branch,draw=msColor1] (4.85,-0.06) to[out=180,in=90] ++(-0.65,-0.48)
    node[below left,ms body,text=msColor1] {biology};
  \draw[branch,draw=msColor1] (5.55,-0.06) to[out=180,in=90] ++(-0.48,-0.48)
    node[below right,ms body,text=msColor1] {environment};
  \draw[branch,draw=msInk] (6.15,0.06) to[out=180,in=-90] ++(-0.60,0.78)
    node[above,ms body] {culture};
  \draw[branch,draw=msColor2] (7.22,0.06) to[out=180,in=-90] ++(-0.45,0.50)
    node[above,ms body] {chronic};
  \draw[branch,draw=msColor2] (8.20,0.06) to[out=180,in=-90] ++(-0.45,0.50)
    node[above,ms body] {acute};
  \draw[branch,draw=msColor6] (8.95,-0.06) to[out=180,in=90] ++(-0.48,-0.60)
    node[below,align=center,ms body,text=msColor6] {releasing\\stimuli};
  \draw[branch,draw=msColor6] (9.42,0.06) to[out=180,in=-90] ++(-0.45,0.60)
    node[above right,align=center,ms body,text=msColor6] {nervous\\system};

  \draw[timespan,draw=msColor5]
    (-0.20,1.04) -- (2.92,1.04)
    node[midway,above=\msSpaceLabel,ms body,text=msColor5] {evolutionary time};
  \draw[timespan,draw=msColor1]
    (6.15,-1.08) -- (3.40,-1.08)
    node[midway,below=\msSpaceLabel,ms body,text=msColor1] {perinatal};
  \draw[timespan,draw=msColor2]
    (6.15,1.04) -- (8.25,1.04)
    node[midway,above=\msSpaceLabel,ms body] {hormones};

  \node[ms small,text=msMuted,anchor=north]
    at (5.35,-2.48) {nested scientific perspectives across timescales};
\end{tikzpicture}
\end{document}
